\documentclass[
	fontsize=12pt,          % default font size 12pt
	paper=a4,               % DIN A4 page format
	numbers=noenddot,       % remove dots behind chapter numbers (e.g. 1.5 not 1.5.)
	listof=totoc,           % add list of figures, tables, etc. to ToC
	listof=entryprefix,     % add entry name to figures, tables, etc.
	listof=nochaptergap,    % no chapter gap for figures, tables, etc.
	bibliography=totoc,     % add bibliography to ToC but without a chapter number
	parskip=half            % half line spacing between paragraphs
	openany                 % chapters can start on any page
]{scrbook}
\usepackage{tcolorbox}

\include{../../../for_latex/preamble}
\title{\Huge\textbf{Längster Stab}}
\author{Luis Staudt}
\date{}

% Code-Highlighting für Java
\definecolor{codegray}{rgb}{0.5,0.5,0.5}
\definecolor{codepurple}{rgb}{0.58,0,0.82}
\definecolor{backcolour}{rgb}{0.95,0.95,0.92}
\definecolor{keywordcolor}{rgb}{0.8,0.47,0.13}

\lstdefinestyle{java}{
	language=,
	basicstyle=\ttfamily\small,
	keywordstyle=\bfseries\color{blue!70!black},
	stringstyle=\color{red!60!black},
	commentstyle=\itshape\color{gray},
	numbers=none,
	numberstyle=\tiny\color{gray},
	numbersep=8pt,
	frame=single,
	framerule=0.4pt,
	rulecolor=\color{gray!50},
	breaklines=true,
	showstringspaces=false,
	tabsize=4,
	xleftmargin=1.5em,
	framexleftmargin=1.5em,
	aboveskip=0.8em,
	belowskip=0.6em
}
\lstset{style=java}

\newtcolorbox{hinweis}[1][Hinweis]{
	colback=blue!5,
	colframe=blue!40!black,
	fonttitle=\bfseries,
	title={#1},
	boxrule=0.5pt,
	arc=2pt
}

\newtcolorbox{beispiel}[1][Beispiel]{
	colback=green!5,
	colframe=green!40!black,
	fonttitle=\bfseries,
	title={#1},
	boxrule=0.5pt,
	arc=2pt
}

\begin{document}

% --- Titelblock ---
\begin{center}
	\rule{\textwidth}{0.8pt}\\[0.6em]
	{\LARGE\bfseries Intensivtutorium}\\[0.3em]
	{\large Programmierung}\\[0.6em]
	\begin{tabular}{r l @{\hspace{2cm}} r l}
		\textbf{Semester:}      & Sommersemester 2026                    & \textbf{Tutor:}         & Luis Staudt \\
		\textbf{Studiengänge:}  & PEB \& MVB         & & \\
	\end{tabular}\\[0.6em]
	\rule{\textwidth}{0.8pt}
\end{center}
\vspace{1em}

% --- Seitenkopf und -fuß (ab Seite 2) ---
\pagestyle{fancy}
\fancyhf{}
\lhead{\small Intensivtutorium, Grundlagen der Programmierung}
\rhead{\small SS 2026}
\cfoot{\thepage}
\renewcommand{\headrulewidth}{0.4pt}


% =============================================
% Aufgabe 4: Längster Stab
% =============================================

\section*{Aufgabe 4: Längster Stab}

Gegeben ist die aus Vorlesung und Übung bekannte Datei \texttt{fachwerk.jar}, die die Funktionalität zur Berechnung von Fachwerken zur Verfügung stellt.
Für eine konstruktive Aufgabe, etwa die Auswahl geeigneter Halbzeuglängen, möchten Sie für ein gewähltes Fachwerk herausfinden, welcher Stab der \textbf{längste} ist.
Die Länge eines Stabes ergibt sich aus den Koordinaten seines Anfangs- und Endknotens über den Satz des Pythagoras:

\[
	l = \sqrt{(x_e - x_a)^2 + (y_e - y_a)^2}
\]

\begin{center}
	\begin{tikzpicture}[>=Stealth, scale=1.0,
		knoten/.style={circle, fill=black, inner sep=1.6pt},
		stab/.style={thick, gray!70!black},
		laengster/.style={very thick, red!75!black}]

		% Knotenkoordinaten
		\coordinate (A) at (0, 0);
		\coordinate (B) at (3, 0);
		\coordinate (C) at (7, 0);
		\coordinate (D) at (1.5, 2.2);
		\coordinate (E) at (5, 2.5);

		% Stäbe
		\draw[stab]      (A) -- (B);
		\draw[laengster] (B) -- (C);   % längster Stab
		\draw[stab]      (A) -- (D);
		\draw[stab]      (D) -- (B);
		\draw[stab]      (D) -- (E);
		\draw[stab]      (E) -- (B);
		\draw[stab]      (E) -- (C);

		% Knoten
		\foreach \p in {A,B,C,D,E} \node[knoten] at (\p) {};
		\node[below left,  font=\footnotesize] at (A) {Knoten 0};
		\node[below,       font=\footnotesize] at (B) {Knoten 1};
		\node[below right, font=\footnotesize] at (C) {Knoten 2};
		\node[above left,  font=\footnotesize] at (D) {Knoten 3};
		\node[above,       font=\footnotesize] at (E) {Knoten 4};

		% Beschriftung längster Stab
		\node[below, font=\footnotesize\color{red!75!black}, yshift=-2pt]
			at ($(B)!0.5!(C)$) {längster Stab};
	\end{tikzpicture}
\end{center}

\medskip
Erweitern Sie den unten abgebildeten Quelltext zu einem vollständigen, lauffähigen Programm.
Die \textbf{Stabnummer} (Nummerierung wie im Programm, beginnend bei~0) des längsten Stabes sowie der \textbf{Wert} seiner Länge sollen dem Anwender auf dem Bildschirm ausgegeben werden.

\begin{hinweis}[Benötigte Methoden aus \texttt{fachwerk.jar}]
	\begin{itemize}[nosep]
		\item \texttt{int gibAnzahlStaebe()}: Anzahl der Stäbe im Fachwerk.
		\item \texttt{Fachwerkstab gibStab(int i)}: der Stab mit Index~\texttt{i}.
		\item \texttt{int gibKnotenNrA()} / \texttt{int gibKnotenNrE()}: Index des Anfangs- bzw.\ Endknotens eines Stabes.
		\item \texttt{Knoten gibKnoten(int i)}: der Knoten mit Index~\texttt{i}.
		\item \texttt{double gibX()} / \texttt{double gibY()}: Koordinaten eines Knotens.
	\end{itemize}
\end{hinweis}

\begin{hinweis}[Wichtig]
	Es geht hier nur um die \emph{Geometrie} des Fachwerks (die Knotenkoordinaten).
	Ein Aufruf von \texttt{berechneGroessen()} ist daher \textbf{nicht} erforderlich.
\end{hinweis}

\bigskip
\textbf{Der abgebildete Quelltext ist als Vorlage gedacht und enthält noch mehrere Fehler} (Tipp\-fehler, falsche Typen, falsche Methodenaufrufe, fehlende Semikolons und einen Logikfehler).
Finden und korrigieren Sie diese, sodass das Programm korrekt kompiliert und das gewünschte Ergebnis liefert.

\begin{lstlisting}
public class Fachwerkaufgabe {
    public static void main(String[] args) {
        Fachwerk fw = null;
        System.out.println("Welches Fachwerk wollen Sie rechnen (0-4)?");
        int wahl = Tastatureingabe.leseInt();
        fw = new Fachwerk(wahl);

        double max = 0;
        int nr = 0;

        for (int i = 0; i < fw.gibAnzahllStaebe(); i++) {
            double a = fw.gibStab(i).gibKnotenNrA();
            double e = fw.gibStab(i).gibKnotenNrE();
            double dx = fw.gibKnoten(e).gibX() - fw.gibKnoten(a).gibX();
            double dy = fw.gibKnoten(e).gibX() - fw.gibKnoten(a).gibY();
            double laenge = Math.sqrt(dx * dx + dy + dy)
            if (laenge > max) {
                max = laenge;
                nr = i
            }
        }
        System.out.pritnln("Laengster Stab: Stab " + Nr + " mit Laenge " + max);
    }
}
\end{lstlisting}

\begin{beispiel}[Mögliche Ausgabe (Fachwerk~2)]
\begin{verbatim}
Welches Fachwerk wollen Sie rechnen (0-4)?
2
Laengster Stab: Stab 7 mit Laenge 4.272001872...
\end{verbatim}
\end{beispiel}

\end{document}
